\begin{figure}[ht]
\centering
\begin{tikzpicture}[x=1cm, y=1cm, font=\small]

% Stage 0: original A
\begin{scope}[shift={(0,0)}]
  \node[anchor=south, font=\small\bfseries] at (1.4, 2.6) {Stage 0};
  \node[anchor=south, font=\scriptsize] at (1.4, 2.2) {original};
  \matrix[matrix of math nodes, left delimiter={(}, right delimiter={)},
          column sep=4pt, row sep=2pt] at (1.4, 1.0) {
    2 & 1 & 1 \\
    4 & 3 & 3 \\
    8 & 7 & 9 \\
  };
\end{scope}

% arrow
\draw[->, thick] (3.2, 1.0) -- (3.8, 1.0) node[midway, above, font=\scriptsize] {pivot $A_{11}$};

% Stage 1
\begin{scope}[shift={(4.0,0)}]
  \node[anchor=south, font=\small\bfseries] at (1.4, 2.6) {Stage 1};
  \node[anchor=south, font=\scriptsize] at (1.4, 2.2) {column 1 zeroed};
  \matrix[matrix of math nodes, left delimiter={(}, right delimiter={)},
          column sep=4pt, row sep=2pt] at (1.4, 1.0) {
    2 & 1 & 1 \\
    |[engblue]| 0 & 1 & 1 \\
    |[engblue]| 0 & 3 & 5 \\
  };
\end{scope}

\draw[->, thick] (7.2, 1.0) -- (7.8, 1.0) node[midway, above, font=\scriptsize] {pivot $A_{22}$};

% Stage 2: U
\begin{scope}[shift={(8.0,0)}]
  \node[anchor=south, font=\small\bfseries] at (1.4, 2.6) {Stage 2};
  \node[anchor=south, font=\scriptsize] at (1.4, 2.2) {$\mat{U}$ ready};
  \matrix[matrix of math nodes, left delimiter={(}, right delimiter={)},
          column sep=4pt, row sep=2pt] at (1.4, 1.0) {
    2 & 1 & 1 \\
    0 & 1 & 1 \\
    0 & |[engblue]| 0 & 2 \\
  };
\end{scope}

% Multipliers stored as L below
\begin{scope}[shift={(2.0, -3.2)}]
  \node[anchor=south, font=\small\bfseries] at (3.0, 1.4) {Multipliers recorded as $\mat{L}$};
  \matrix[matrix of math nodes, left delimiter={(}, right delimiter={)},
          column sep=4pt, row sep=2pt] at (3.0, 0.2) {
    1 & 0 & 0 \\
    |[engred]| 2 & 1 & 0 \\
    |[engred]| 4 & |[engred]| 3 & 1 \\
  };
  \node[anchor=west, font=\scriptsize, text width=4.6cm] at (4.8, 0.2) {
    Row 2: $R_{2} - 2 R_{1}$. \\
    Row 3: $R_{3} - 4 R_{1}$, then $R_{3} - 3 R_{2}$. \\
    Diagonal of $\mat{L}$ is unit; entries below diagonal store the multipliers.
  };
\end{scope}

\end{tikzpicture}
\caption[Two-stage Gaussian elimination on a $3 \times 3$ system]{Gaussian
elimination on the worked example of section 9.3. Stage 1 uses the
pivot $A_{11} = 2$ to zero entries below; stage 2 uses the new pivot
$A_{22} = 1$ to zero the entry below it. The resulting upper-triangular
matrix is $\mat{U}$; the row multipliers used at each step, recorded in
the lower triangle with unit diagonal, give $\mat{L}$, and
$\mat{L}\mat{U} = \mat{A}$. For the back substitution the entries of
$\vect{b}$ are transformed in lockstep.}
\label{fig:vol02:ch09:gaussian-elimination}
\end{figure}
